%Befehl um Einheiten im Acronym ausgeben zu können
\newcommand{\acrosecondcolumn}[1]{\acroextra{\makebox[21mm][l]{#1}}}

% ============= Eingabe der Symbole und Abkürzungen ============= 
\chapter*{Abkürzungsverzeichnis}\thispagestyle{rmplain}
\addcontentsline{toc}{chapter}{Abkürzungsverzeichnis}
\label{sec:abkuerzungsverzeichnis}
\begin{acronym}[LONGEST]
	\acro{KDE}{K Desktop Environment}
	\acro{SQL}{Structured Query Language}
	\acro{Bash}{Bourne-again shell}
	\acro{JDK}{Java Development Kit}
	\acro{VM}{Virtuelle Maschine}
	%Falls eine Abkürzung in der Mehrzahl nicht einfach auf "s" endet muss das speziell eingestellt werden. 
	\acro{slmtA}{super lange mega tolle Abkürzung} %Einzahl
	\acroplural{slmtA}[slmtAs]{super lange mega tolle Abkürzungen} %Mehrzahl
\end{acronym}


\chapter*{Symbolverzeichnis}\thispagestyle{rmplain}
\addcontentsline{toc}{chapter}{Symbolverzeichnis}
\label{sec:symbolverzeichnis}
\begin{acronym}[LONGEST]
	\acro{temp}[$^\circ\text{C}$]{Europäische Einheit der Temperatur}
	\acro{mu}[$\mu$]{Reibungskoeffizient}
	\acro{pi}[$\pi$]{Kreiszahl}
	\acro{alpha}[$\alpha$]{Winkel zwischen $0^\circ\text{C}$ und $360^\circ\text{C}$}
	\acro{n}[$\mathbb{N}$]{Menge der natürlichen ganzen Zahlen}
	\acro{A}[\ensuremath{A}]{Oberfläche}
	%\acro{A}[\ensuremath{A}]{\acrosecondcolumn{$m^2$}Oberfläche}	%Darstellung mit Einheit
\end{acronym}